\section{The Misunderstanding of Paganism}

\begin{quotex}
In order to illustrate some of the points made recently, we have invited an old friend, Tony Ciapo, to provide translations, with commentary, of some of Evola's works from \textit{Sintesi di dottrina della razza}, \textit{Defesa della razza}, and Evola's commentary on the Fascist movements of the first half of the 20\textsuperscript{th} century. By exploring certain movements from the perspective of Tradition---which are far from any modernist misunderstandings---, we expect to point who are in \textbf{Solidarity} with the few remaining men of Tradition left today and which viewpoints are in \textbf{Continuity} with Tradition itself. Tony will continue this task as long as interest is expressed and intelligent discussion ensues. The initial section is the fourth in the chapter ``La razza aria e il problema spirituale''. His commentaries will appear inside boxes or between square brackets. There is no implication that either Evola or myself would have accepted all such.
\end{quotex}

\paragraph{The Misunderstanding of the new racist Paganism}


Having made the problem clear in such a way, it is perhaps appropriate to point out the misunderstanding---a misunderstanding of no small moment--- characteristic of those contemporary extremist racist currents that believe that have resolved it in terms of a neopaganism. Such a misunderstanding, in truth, is already apparent precisely in the use of terms such as ``pagan'' and ``paganism''. We ourselves had at one point adopted them and sincerely regret it.

\begin{quotex}
This section follows the section ``Ex Occidente Lux --- The Religious Problem'' (to be translated next), which is alluded to in the first sentence. There he discusses Catholicism, in particular its teaching of a primordial revelation of the Fathers. Attentive readers of Gornahoor will recall that Catholicism regards itself as the latest manifestation of the primordial religion. Evola digs deeper to the pre-Christian religion of the \textbf{Ancient City}, whose lineaments we have been at pains to point out. In order to avoid the ambiguities and attendant misunderstandings of the term ``paganism'', we will henceforth refer to it as the ``\textbf{Old Religion}``. It should be obvious that the ``contemporary currents'' he opposes include the amalgam of racism and neopaganism associated with the theoreticians of National Socialism. There is no need to name names, as Evola means all of them. The book he sincerely regrets is Pagan Imperialism.
\end{quotex}

Certainly the word, \emph{paganus}, appears in ancient Latin writers such as Livy with no special purpose. But this does not alter the fact that with the arrival of the new faith, the word \emph{paganus} took on a decidedly pejorative meaning, adopted for use in early Christian apologetics. \emph{Paganus} derives from \emph{pagus}, meaning a village or slum, so that \emph{paganus} means what is characteristic of a country hick or an uncultured, primitive man. In order to profess and glorify the new faith, a certain Christian apologistics, following the the vice of overvaluing themselves through the discrediting of others, proceeded to a systematic and often conscious disparagement and misrepresentation of almost all the earlier doctrines, forms of worship, and traditions, which were made to correspond to the combined and pejorative designation ``paganism''. Naturally, with this goal, it took care to highlight everything that, in the pre-Christian religions and traditions, lacked any normal or primordial character, rather than the clear significance of degenerate and decadent forms. Such a polemical \emph{animus} then led, in particular, to attributing indiscriminately an anti-Christian character to everything that, prior to Christianity, could also simply be non-Christian and could not constitute an irreducible antithesis.

\begin{quotationx}
Evola is making the point that the Old Religion, much older than Christianity, was not, and could not be, explicitly anti-Christian. Therefore, we have to distinguish the former from the paganism prevalent at the beginning of Christianity. By then the Old Religion had been largely forgotten and even when its rites and traditions were followed, it was done with little understanding. The Roman priesthood had been reduced to a formalism, while the aristocrats were more likely to adopt one of the many prevalent Greek philosophies or dabble in the so-called mysteries. As was recently pointed out, it was the plebs, initially without a religion, who were the ``pagans'' of the era. The first Christians came either from among the Hellenized Jews or the Romans who had converted to Judaism whether formally or by adopting their beliefs and ways. (Note that the Judaism at the time was quite different from the contemporary religion that goes by that name.) These would have been predominantly the third caste working as artisans, traders, and other pecuniary occupations.

Of course, the neopagans, following Nietzsche, now return the favor, by claiming it was really the Christians who were the ``hicks''. Such mutual name-calling is unhelpful and pointless. Instead, \emph{pace} the early apologists and neopagans, we follow \textbf{Valentin Tomberg} `s Hermetic maxim, \emph{we must love our pagan past}, provided we understand by that the Old Religion.
\end{quotationx}

On such a basis, it is therefore necessary to think that there is a ``paganism'', meaning something essentially and tendentiously constructed. That is, it lacks any true correspondence to historical reality. We mean, correspondence to what the pre-Christian and especially Aryan world always was in its ``normal'' forms, and not only in decadent aspects or that were refashioned from the degenerate residues of older civilizations or inferior races.

\begin{quotex}
Recommended background material:

\textbf{Fustel}'s \emph{The Ancient City} is, by his own admission, indispensable to understanding what Evola means by the pre-Christian Tradition.

For the history of the Early Christians and Jews, I rely upon the works of \textbf{Ernest Renan}.
\end{quotex}

\paragraph{Neopaganism and Naturalism}

Keeping all this in mind, we come today to discover a peculiar paradox: to wit, such a paganism never even existed but was constructed by Christian apologetics. Yet some ``pagan'' and anti-Christian movements of racism and extreme nationalism very frequently accept it, thus threatening to make it become true today for the first time in history. No more and no less than that.

\begin{quotex}
This critique is devastating. Instead of recovering the primordial religion of their ancestors, these neopagans are adopting a caricature of that religion. It is obvious that the primary characteristics of neopaganism (anti-Christian, racism, nationalism) had no analog in ancient times and are purely anachronisms.
\end{quotex}

What are the principle traits of the pagan outlook on life, as its apologists believe and popularize?

\paragraph{Naturalism}

First and foremost: \textbf{naturalism}. All transcendence is totally unknown to the pagan view of life. It would remain in a mixture of spirit and nature, in an ambiguous unity of body and soul. There is nothing to its religion but a superstitious deification of natural phenomena or the energies of the races, elevated as so many idols. From this, in the first place, a particularism, a polytheism conditioned by blood and soil. In the second place, an absence of the spiritual personality and freedom, a state of innocence that is simply characteristic of beings of nature, of those not yet awakened to any truly supernatural aspiration. This ``innocence'': license, ``sin'', and the delights of the flesh. Even in other domains, either superstition, or a purely ``profane'', materialistic, and fatalistic civilization. Apart from certain ``anticipations'' considered insignificant, it is with Christianity that the world of supernatural would be revealed for the first time, that is, the world of grace and personality, in contrast to ``pagan'' determinism and naturalism. With it a ``catholic'' ideal (in the etymological sense of universality) would be affirmed; a healthy dualism, permitting the subordination of nature to a higher law, and the triumph of the law of the spirit beyond the law of flesh, blood, and the false gods.

\begin{quotex}
Here Evola identifies Naturalism as the primary trait of the neopagans. This is true for the most crude attempts at revival such as Wicca or Asatru, as well as the more sophisticated systems. To deny this, and thus to affirm supernaturalism, freedom, and personality, would draw neopaganism closer to Christianity, at least in a common metaphysical perspective. We should point out that Evola opposes Christianity on completely different grounds from those of the neopagans, who simply repeat the slogans of the Pharisees and the Enlightenment, sworn enemies of Tradition.
\end{quotex}

These are the most typical traits of the predominant conception of paganism, i.e., of everything that is not specifically the Christian worldview. Anyone who possesses any direct acquaintance with cultural and religious history, however elementary, can see how inaccurate and one-sided this attitude is. Besides, some early Church Fathers often demonstrated a greater understanding of the symbols and cults of preceding civilizations. Here we highlight just a few points.

\begin{quotex}
Here Evola's frustration is showing: even an elementary knowledge of history and religion would suffice to refute the bases of neopaganism, yet we cannot even count on that today. The Church Fathers showed a deeper understanding of the Old Religion than do our own neopagans. Regular readers of Gornahoor will find many references to this.
\end{quotex}


\flrightit{Posted on 2011-07-17 by Aeneas}

\begin{center}* * *\end{center}

\begin{footnotesize}
\begin{sffamily}
\texttt{Perennial on 2011-07-18 at 01:17 said:}

Great stuff! Look forward to reading subsequent installments.

\hfill

\texttt{James O'Meara on 2011-07-19 at 23:42 said:}

So, Evola rejected the crude ``paganism'' of an Alfred Rosenberg. Apart from already knowing that, what does all this effort provide me with? Even Hitler agreed ``When the Greeks built temples, we were living in mud huts. Why does Himmler insist on digging all this up?'' Evola collaborated with Nazi pagans, which the modern PC types will never accept, no matter what ``distinctions'' you can find. As far as `racism' and `true Christianity' go, the bottom line is that Evola would have sent M L King to a death camp without a second thought, so what have you proved about his ``sophisticated'' versions of racism, paganism and state power?

AS for ``Evola regretted his book Pagan Imperialism,'' well, yes, he was unable to influence Mussolini so he turned to the Nazis, hoping to ``rectify'' Path of Cinnabar their crude racism and paganism. Again, Rachel Maddow isn't going to `friend' him, so what's the point?

As for your sneers I won't call it a `critique' against Nietzsche, Fred spoke about `intellectual conscience.' Anyone with an intellectual conscience would point out that while Evola may have `regretted' Pagan Imperialism, his last book, Ride the Tiger, is based on the idea that all honest modern thinkers `intellectual conscience' again will take Nietzsche as their starting point. Since he lists Nietzsche as his most profound teenage influence, we can see that Evola was consistently Nietzschean from age 16 to his death. Now there's integrity!

BTW, my own reflections on Nietzsche, the Greeks and the function of religion can be found here:

\url{http://tinyurl.com/3glcsx9}


\hfill

\texttt{Cologero on 2011-07-20 at 00:02 said:}

It takes considerable effort, James, to provide original translations and context to the masses. You could take the high road and thank us.

\hfill

\texttt{Perennial on 2011-07-20 at 00:36 said:}

Goodness James, bit negative, are we not?

\hfill

\texttt{Aeneas on 2011-07-20 at 08:19 said:}

Consistently a Nietzschean? Pas de tout! Evola writes in Cinnabar regarding the early influence of Nietzsche:

``He fed a basic tendency, even if in confused and partly distorted forms, therefore a mixture of positive and negative.''

Evola also refers to the ``ambiguous character of many view of Nietzsche''.

Regarding Nietzsche's personal life, Evola attributed it to the ``awakeing of the elemnt of transcendence without its conscious and active assuption.'' Evola refers to the ideal of the Superman as the ``worst of Nietzsche.''

So that is the worst of Evola, that he tries to be ``hip'' with references to Nietzsche, while admitting what a mixed bag he really is. Evola's position was to critique the false ideals of his time and Nietzsche served that purpose. That task having been accomplished, we at Gornahoor are delineating the characteristics a the Tradition to follow, where there is no need or room for Nietzsche. As for ``intellectual integrity'', a ``scholar'' needs to take into account the entire corpus of an author. We have access to thousands of pages of Evola that few ``Evolians'' bother to familiarize themselves with. Hence, we try to make them available via translations and commentaries. Those with intellectual integrity will appreciate these efforts. As for the others, perhaps we can straighten them out.

\hfill

\texttt{I on 2011-07-20 at 08:54 said:}

``As far as `racism' and `true Christianity' go, the bottom line is that Evola would have sent M L King to a death camp without a second thought, so what have you proved about his ``sophisticated'' versions of racism, paganism and state power?''

James, this is not true. Evola was firmly and consistently against violent treatment of political enemies and ideological opponents. This is something that made him un-popular among some high ranking fascists and nazis of the social darwinist type, who believed in the ``Nietzshean'' myths of the blond beast of prey etc. It is true that he would have made and accepted all serious efforts to remove the kinds of ML King from places of influence and (political) power, but he would not have suggested a death camp or execution for them. This is not because he had felt any humanistic empathy for them, but he could see where such actions would inevitably and objectively lead -- to disaster and ruin.

\hfill

\texttt{Max on 2015-04-19 at 06:24 said:}

Tradition is lived in a specific time and place, but some treat it like existing in a vacuum outside real societies, humans and events, without which it becomes just another mental construct blocking our view. This is the instance of a mind that only goes backwards but not forward. A fault in the ``thinking'' of modern people is that they start with the effect, and then try to work out which cause lead to it which means they have to make it up, rather than to start at the beginning and align ourselves with that which is given.

\hfill

\texttt{Visionsofglory14 on 2011-07-22 at 04:13 said:}

Perhaps he was referring to your claim that Neo-pagan critiques of Christianity parallel those of the Pharisees and the Enlightenment thinkers. Care to extrapolate?

\hfill

\texttt{Cologero on 2011-07-22 at 12:48 said:}

Mr. Vision, is there some specific point you want to address? First of all, the neo-pagan new right, in its criticism, aligns itself with those said forces of revolution; thus it is not traditional, and not even of the right. I translated Evola's objections to Catholicism (what he and we mean by Christianity) and it is available on this web site. Unlike the crude critiques that are all too available on the www, Evola accepts the possibility of miracles, defends the concept of infallibility, and even the idea of a High Priest (or pope). Furthermore, unlike the neo-pagan new right, Evola defends Charlemagne and regards the Middle Ages as an age of tradition.

I sent out to the mailing list the next installment of the translation of Sintesi, where Evola writes:

\begin{quote}\footnotesize
Once all that has been established, the suggested possibility of ``transcending'' certain aspects of Christianity would become real. According to its strict Latin etymology, to transcend means to surpass while ascending. In principle, it is not a matter -- it is good to repeat it -- of rejecting Christianity or of showing the same incomprehension towards it as Christianity itself has shown, and largely continues to show, towards paganism. It would instead perhaps be a matter of integrating it into something much larger, omitting some aspects, for which it is little in accord with the spirit typical of some of the renovating forces of today, of the transalpine type, in order to accentuate instead other more essential aspects which such faith cannot contradict the general conceptions of Aryan and Nordic, pre-Christian and non-Christian spirituality.
\end{quote}


Has anyone in the neo-pagan new right mentioned the idea of ``transcending'' Christianity?

\hfill

\texttt{Cologero on 2011-07-23 at 00:07 said:}

How many times do we have to make myself clear, Eddie? Just one more time. Point 1. Evola was not at all pro-Christian ... he could not relate to its theology at all. But neither was he anti-Christian. Rather, he was non-Christian, as though it never existed for him. Point 2. When we use the word ``Christian'' at Gornahoor are referring specifically to the Church of the Middle Ages, especially when it was undivided. By that standard, there are precious few ``Christians'' around today. Let's not be confused about labels; we're trying to penetrate deeper. Even failure is worth the attempt. That is one reason that we don't write much about theology proper, which is, as you say, in its contemporary manifestations to a large extent sentimental. I can well understand that it does not appeal to many today. Nevertheless, we need to understand why it did appeal to many a millennium ago, and was the guiding spirit to the last traditional civilization in the West.

A corollary is that we try to transcend the ``us vs them'' attitude, that is all too common among the neo-pagans. Any careful reader can determine that we are not distorting Evola. However, it seems to us that the neo-pagans distort him while claiming him as one of their own. I'll summarize what I see are the specific differences. If there are any wise pagans who address these issues, please point them out to us.

The real paganism is that of the Ancient City about which Gornahoor has devoted several posts. It is not the decadent form of the late Empire. The real paganism was theocratic and was led by a priest-king until its subsequent degenerations. All of life was regulated by duties and rites.

The real paganism was not a hedonistic, sensualistic, naturalism.

The Christian kings and emperors of the Middle Ages, such as Arthur, Charlemagne, Frederick II were exercised legitimate authority.

The Church of the Middle Ages (and the civilization that grew up alongside it) was not Jewish creation initiated by ``Saul''. It was the creation of the Nordic-Roman spirit.

The end of Tradition in the West did not begin with Constantine. It began 1000 years later with the Renaissance, Reformation, Enlightenment, French Revolution.

The Church of the Middle Ages was hierarchical, not in the least democratic and egalitarian.

Those are some points. Evola's view is the first of each item. Every neo-pagan I've read has adopted the second. So, this is not meant to be a ``court case'' between Christians and pagans, if only for the reason that few Christians today would even recognize themselves in this characterization. The issue is to determine the spiritual forces at work and to understand the flow of history apart from angry and ignorant polemics.

\hfill

\texttt{Perennial on 2011-07-23 at 02:38 said:}

Bravo!

\hfill

\texttt{Eumaeus on 2011-07-26 at 16:42 said:}

New to this forum. Excellent conversation.

Indeed many of the early Church Fathers did have a superior understanding of the Old Religion in part because they practiced it prior to conversion. The process of integration of the new religious idiom with the old was an important part of the project of the Fathers although the work was not acknowledged as such. In Exoteric religion they generally do not give much attribution to where they have taken their ideas and indeed the official stories may conceal deeper connections outside of historical traditions. For example the Zoroastrian and Hellenistic inputs into Christianity even at the level of the Gospels are obvious but not footnoted. They are acknowledged by allusion-- the appearance of the Magi at Bethlehem; references to the Lake of Fire-- and the very notion of incarnation already so well recounted for many Greek heroes. And at the core: equation of Christ as Logos.

These are also symbols and markers within the Gospel which betoken the expansion of Christian mission from tribal to universal-- by reference to those two traditions that were already beyond tribal particularism at that point in time-- Zoroastrianism and Hellenism.


\hfill


\end{sffamily}
\end{footnotesize}

